\subsection{补充说明}\label{sec:17f}
如前可证明，定义 $(d,l)$ 时所涉及的选择实质上无关紧要。更准确地，若 $(d,l)$ 与 $(\tilde d,\tilde l)$
\SourcePage{257}
由不同选择定义，则可考虑在两种构造版本之间插值的附加模空间，从而得到
\[
\boxed{\begin{gathered}
\delta:CF^{\mathrm{pr}}(Q,L_1)\longrightarrow CF^{\mathrm{pr}}(Q,L_2),\\
\lambda:CF^{\mathrm{pr}}(L_0,L_1)\otimes CF^{\mathrm{pr}}(Q,L_0)\longrightarrow CF^{\mathrm{pr}}(Q,L_2),\\
\mu^1(\delta(\cdot))+\delta(\mu^1(\cdot))+d+\tilde d=0,\\
\delta(\mu^2(\cdot,\cdot))+\mu^1(\lambda(\cdot,\cdot))
+\lambda(\mu^1(\cdot),\cdot)\\
{}+\lambda(\cdot,\mu^1(\cdot))+l+\tilde l=0.
\end{gathered}}
\]
其代数含义是，在链复形\eref{17.16}之间得到同构
\begin{equation}\label{eq:17.17}
\begin{pmatrix}
\Id\otimes\Id&0&0\\0&\Id&0\\\lambda&\delta&\Id
\end{pmatrix}:D\longrightarrow\widetilde D.
\end{equation}
更一般地，可考察 $(d,l)$ 如何依赖于 $L_1$。假设有与 $L_1$ 同痕的另一个子流形 $\widetilde L_1$，以及相应版本 $(\tilde d,\tilde l)$。取
\SourcePage{258}
延拓映射
\[
\begin{aligned}
a &:CF^{\mathrm{pr}}(Q,L_1)\longrightarrow CF^{\mathrm{pr}}(Q,\widetilde L_1),\\
b &:CF^{\mathrm{pr}}(L_0,L_1)\longrightarrow CF^{\mathrm{pr}}(L_0,\widetilde L_1);
\end{aligned}
\]
它们在链同伦意义下与三角形乘积相容，即存在 $h:CF^{\mathrm{pr}}(L_0,L_1)\otimes CF^{\mathrm{pr}}(Q,L_0)\to CF^{\mathrm{pr}}(Q,\widetilde L_1)$，使
\[
\begin{aligned}
\mu^1(h(\cdot,\cdot))+h(\mu^1(\cdot),\cdot)+h(\cdot,\mu^1(\cdot))\\
{}+\mu^2(b(\cdot),\cdot)+a(\mu^2(\cdot,\cdot))=0.
\end{aligned}
\]
从一般观点看，延拓映射来自平凡纤维化 $M\times Z\to Z$，其中边界条件在 $\R\times\{0\}$ 上恒定等于 $Q$ 或 $L_0$，在 $\R\times\{1\}$ 上则实现从 $\widetilde L_1$ 到 $L_1$ 的同痕；\secref{8j}中已指出这一点。可取三穿孔圆盘上的平凡纤维化，带边界条件 $(Q,L_0,L_1)$，并将其入向端与定义 $a$ 的纤维化的出向端粘合。另一方面，可对三穿孔圆盘取边界条件 $(Q,L_0,\widetilde L_1)$，将第二个出向端与定义 $b$ 的纤维化的入向端粘合。如此构造的两个纤维化可相互变形，所得族给出 $h$。

同样的想法可如下联系 $d$ 与 $\tilde d$：将定义 $a$ 的纤维化的入向端，与 $E^{\tilde d}\to S^{\tilde d}$ 的出向端粘合，再将结果变形回原来的 $E^d\to S^d$。结果是单参数纤维化族 $E^{\alpha,r}\to S^{\alpha,r}$，如\sr{fig:17.6}{图~17.6}；其相应不变量具有形式
\[
\boxed{\begin{gathered}
\alpha:CF^{\mathrm{pr}}(Q,L_1)\longrightarrow CF^{\mathrm{pr}}(Q,L_2),\\
\mu^1(\alpha(\cdot))+\alpha(\mu^1(\cdot))+\tilde d(a(\cdot))+d(\cdot)=0.
\end{gathered}}
\]
\SourceFigure{figure-17-6}{0.62515317}{17.6}

\SourcePage{259}
同类而稍复杂的论证也适用于 $l,\tilde l$，结果具有形式
\[
\boxed{\begin{gathered}
\beta:CF^{\mathrm{pr}}(L_0,L_1)\otimes CF^{\mathrm{pr}}(Q,L_0)\longrightarrow CF^{\mathrm{pr}}(Q,L_2),\\
\mu^1(\beta(\cdot,\cdot))+\beta(\mu^1(\cdot),\cdot)+\beta(\cdot,\mu^1(\cdot))\\
{}+\tilde l(b(\cdot),\cdot)+l(\cdot,\cdot)+\alpha(\mu^2(\cdot,\cdot))
+\tilde d(h(\cdot,\cdot))=0.
\end{gathered}}
\]
合在一起，便得到链映射
\begin{equation}\label{eq:17.18}
\begin{pmatrix}
b\otimes\Id&0&0\\h&a&0\\\beta&\alpha&\Id
\end{pmatrix}:D\longrightarrow\widetilde D.
\end{equation}
\eref{17.18}中的对角项是延拓映射，因而为拟同构，所以整个映射也为拟同构。最后，类似论证表明，$D$ 在 $L_2$ 的同痕下于拟同构意义内保持不变；为避免重复，省略这一部分。
